% ==============================================================================
% SECTION 04.01: PROBABILITY DENSITY FUNCTION (PDF) AND CONTINUOUS SUPPORT
% ==============================================================================

\subsection*{Problem Statements --- Section 04.01}

\subsubsection*{Fundamental Level}

\begin{problema}
	\label{prob:4.1.1}
	Consider the function $f(x) = c \cdot e^{-x}$ for $x \ge 0$ and $f(x) = 0$ for $x < 0$, where $c$ is a positive constant.
	\begin{enumerate}
		\item Determine the value of $c$ such that $f(x)$ is a valid probability density function (PDF).
		\item Compute the probability $P(1 \le X \le 3)$ by integrating the PDF over the interval.
		\item Compute the cumulative distribution function (CDF) $F(x)$ for all $x \in \mathbb{R}$.
	\end{enumerate}
\end{problema}

\begin{problema}
	\label{prob:4.1.2}
	A continuous random variable $X$ has PDF given by:
	$$f(x) = \begin{cases} c \cdot (4x - 2x^{2}) & \text{if } 0 \le x \le 2 \\ 0 & \text{otherwise} \end{cases}$$
	\begin{enumerate}
		\item Find the normalization constant $c$ such that $\int_{-\infty}^{\infty} f(x)\,dx = 1$.
		\item Compute $P(X \ge 1)$ and $P(0.5 \le X \le 1.5)$.
	\end{enumerate}
\end{problema}

\begin{problema}
	\label{prob:4.1.3}
	Let $X$ be a random variable with triangular PDF:
	$$f(x) = \begin{cases} x & \text{if } 0 \le x \le 1 \\ 2 - x & \text{if } 1 < x \le 2 \\ 0 & \text{otherwise} \end{cases}$$
	\begin{enumerate}
		\item Verify that $\int_{-\infty}^{\infty} f(x)\,dx = 1$.
		\item Compute the mode (the value where $f$ is maximal) and the median (the value $m$ such that $P(X \le m) = 0.5$).
	\end{enumerate}
\end{problema}

\subsubsection*{Operational Level}

\begin{problema}
	\label{prob:4.1.4}
	Let $X$ be a continuous random variable with CDF $F(x) = 1 - e^{-x^{2}/2}$ for $x \ge 0$ and $F(x) = 0$ for $x < 0$ (Rayleigh distribution).
	\begin{enumerate}
		\item Derive the PDF $f(x)$ by applying $f(x) = F'(x)$.
		\item Verify that $\int_{0}^{\infty} f(x)\,dx = 1$ using the substitution $u = x^{2}/2$.
		\item Compute the median $m$ such that $F(m) = 0.5$.
	\end{enumerate}
\end{problema}

\begin{problema}
	\label{prob:4.1.5}
	Consider the PDF defined by:
	$$f(x) = \begin{cases} \dfrac{3}{8}(2 - x^{2}) & \text{if } -1 \le x \le 1 \\ 0 & \text{otherwise} \end{cases}$$
	\begin{enumerate}
		\item Verify that $f$ integrates to $1$.
		\item Compute $P(-0.5 \le X \le 0.5)$ and $P(X \ge 0)$.
		\item Find the CDF $F(x)$ for $x \in [-1, 1]$ and qualitatively sketch its shape.
	\end{enumerate}
\end{problema}

\begin{problema}
	\label{prob:4.1.6}
	In a communications system, the waiting time $X$ (in milliseconds) until a packet is received has an exponential PDF with parameter $\lambda = 0.5$ ms$^{-1}$:
	$$f(x) = 0.5 \cdot e^{-0.5 x}, \quad x \ge 0$$
	\begin{enumerate}
		\item Compute the probability $P(X \le 1)$ that the waiting time is less than 1 ms.
		\item Compute $P(2 \le X \le 5)$ by integrating the PDF.
		\item Find the median $m$ such that $P(X \le m) = 0.5$.
	\end{enumerate}
\end{problema}

\subsubsection*{Analytical Level}

\begin{problema}
	\label{prob:4.1.7}
	Let $X$ be a continuous random variable with PDF $f(x)$. Formally prove that the mathematical expectation can be computed as:
	$$\E[X] = \int_{-\infty}^{\infty} x \cdot f(x)\,dx$$
	starting from the definition $\E[X] = \int x\,dF(x)$ and using integration by parts under technical regularity conditions.
\end{problema}

\begin{problema}
	\label{prob:4.1.8}
	For a continuous random variable $X$ with PDF $f(x)$ and a function $g: \mathbb{R} \to \mathbb{R}$, prove the \emph{Continuous LOTUS Theorem} (Law of the Unconscious Statistician):
	$$\E[g(X)] = \int_{-\infty}^{\infty} g(x) \cdot f(x)\,dx$$
	\begin{enumerate}
		\item Prove the particular case $g(X) = X^{2}$.
		\item Derive the variance formula $\Var(X) = \E[X^{2}] - (\E[X])^{2}$ using LOTUS with $g(x) = (x - \mu)^{2}$.
	\end{enumerate}
\end{problema}

\subsubsection*{Challenging Level}

\begin{problema}
	\label{prob:4.1.9}
	Consider a mixed continuous-discrete distribution where $X$ takes the value $0$ with probability $1/3$ and, conditionally on $X > 0$, follows an exponential distribution with parameter $\lambda = 2$:
	$$f(x) = \begin{cases} 1/3 & \text{if } x = 0 \\ (2/3) \cdot 2 e^{-2x} & \text{if } x > 0 \\ 0 & \text{otherwise} \end{cases}$$
	\begin{enumerate}
		\item Verify that $f$ is a valid PDF (point mass at $0$ + exponential density).
		\item Compute the CDF $F(x)$ for all $x \in \mathbb{R}$, identifying the jump at $x = 0$.
		\item Compute $\E[X]$ and $\Var(X)$ using the law of total expectation.
	\end{enumerate}
\end{problema}

\begin{problema}
	\label{prob:4.1.10}
	In signal theory, the convolution of two densities $f$ and $g$ produces the density of the sum $Z = X + Y$ of two independent random variables:
	$$h(z) = (f * g)(z) = \int_{-\infty}^{\infty} f(x) \cdot g(z - x)\,dx$$
	\begin{enumerate}
		\item Prove that if $X \sim N(0, 1)$ and $Y \sim N(0, 1)$ are independent, then $Z = X + Y$ has PDF $N(0, 2)$, that is, $\Var(Z) = \Var(X) + \Var(Y)$.
		\item Generalize: if $X_{1}, \dots, X_{n}$ are i.i.d. $N(0, \sigma^{2})$, then $\sum_{i=1}^{n} X_{i} \sim N(0, n\sigma^{2})$.
		\item Interpret this result in the context of Gaussian noise filtering in telecommunications.
	\end{enumerate}
\end{problema}

\subsection*{Hints for the Problems --- Section 04.01}

\begin{sugerencia}
	\label{sug:4.1.1}
	For Problem \ref{prob:4.1.1}, integrate $\int_{0}^{\infty} c \cdot e^{-x}\,dx = c = 1$, so $c = 1$. For $P(1 \le X \le 3) = e^{-1} - e^{-3} \approx 0.318$. The CDF is $F(x) = 1 - e^{-x}$ for $x \ge 0$.
\end{sugerencia}

\begin{sugerencia}
	\label{sug:4.1.2}
	For Problem \ref{prob:4.1.2}, integrate $\int_{0}^{2} c(4x - 2x^{2})\,dx = c \cdot [2x^{2} - 2x^{3}/3]_{0}^{2} = c(8 - 16/3) = 8c/3 = 1$, giving $c = 3/8$. For $P(X \ge 1)$, integrate from $1$ to $2$.
\end{sugerencia}

\begin{sugerencia}
	\label{sug:4.1.3}
	For Problem \ref{prob:4.1.3}, the mode is at $x = 1$ (where the PDF reaches its maximum of $1$). The integral $\int_{0}^{1} x\,dx = 1/2$ and $\int_{1}^{2} (2-x)\,dx = 1/2$, summing to $1$. The median $m$ satisfies $\int_{0}^{m} x\,dx = 0.5$ if $m \le 1$, giving $m = 1$.
\end{sugerencia}

\begin{sugerencia}
	\label{sug:4.1.4}
	For Problem \ref{prob:4.1.4}, $f(x) = F'(x) = x e^{-x^{2}/2}$ for $x \ge 0$. With the substitution $u = x^{2}/2$, $du = x\,dx$: $\int_{0}^{\infty} x e^{-x^{2}/2}\,dx = \int_{0}^{\infty} e^{-u}\,du = 1$. The median satisfies $1 - e^{-m^{2}/2} = 0.5$, giving $m = \sqrt{2 \ln 2} \approx 1.177$.
\end{sugerencia}

\begin{sugerencia}
	\label{sug:4.1.5}
	For Problem \ref{prob:4.1.5}, $\int_{-1}^{1} (3/8)(2 - x^{2})\,dx = (3/8)[2x - x^{3}/3]_{-1}^{1} = (3/8)(4 - 2/3) = (3/8)(10/3) = 10/8 \ne 1$. Check: $\int_{-1}^{1} (2 - x^{2})\,dx = [2x - x^{3}/3]_{-1}^{1} = (2 - 1/3) - (-2 + 1/3) = 4 - 2/3 = 10/3$. Therefore $c = 3/10$ for the integral to equal $1$.
\end{sugerencia}

\begin{sugerencia}
	\label{sug:4.1.6}
	For Problem \ref{prob:4.1.6}, the exponential CDF is $F(x) = 1 - e^{-0.5 x}$. Then $P(X \le 1) = 1 - e^{-0.5} \approx 0.393$. $P(2 \le X \le 5) = e^{-1} - e^{-2.5} \approx 0.3679 - 0.0821 = 0.286$. The median is $m = -\ln(0.5)/0.5 = 2 \ln 2 \approx 1.386$ ms.
\end{sugerencia}

\begin{sugerencia}
	\label{sug:4.1.7}
	For Problem \ref{prob:4.1.7}, use the representation $F(x) = \int_{-\infty}^{x} f(t)\,dt$ and integration by parts $\int x\,dF(x) = xF(x) - \int F(x)\,dx$. Assuming $\lim_{x \to \pm\infty} xF(x) = 0$ (technical light-tail condition), substitute the CDF and simplify to obtain $\E[X] = \int x f(x)\,dx$.
\end{sugerencia}

\begin{sugerencia}
	\label{sug:4.1.8}
	For Problem \ref{prob:4.1.8}, use the fact that for continuous random variables, $\E[g(X)] = \int g(x) f(x)\,dx$ where $f$ is the density. The case $g(X) = X^{2}$ follows directly. For the variance, $g(X) = (X - \mu)^{2}$ with $\mu = \E[X]$, giving $\E[(X-\mu)^{2}] = \E[X^{2}] - 2\mu\E[X] + \mu^{2} = \E[X^{2}] - \mu^{2}$.
\end{sugerencia}

\begin{sugerencia}
	\label{sug:4.1.9}
	For Problem \ref{prob:4.1.9}, the mass at $0$ is $1/3$ and the rescaled exponential density $(2/3) \cdot 2 e^{-2x}$ integrates to $\int_{0}^{\infty} (4/3) e^{-2x}\,dx = (4/3) \cdot (1/2) = 2/3$. Summing: $1/3 + 2/3 = 1$. The CDF has a jump of $1/3$ at $x = 0$ and then increases according to the rescaled exponential.
\end{sugerencia}

\begin{sugerencia}
	\label{sug:4.1.10}
	For Problem \ref{prob:4.1.10}, the convolution of two gaussians is gaussian with summed variances. Prove this using the explicit form of the PDF: $h(z) = \int (2\pi)^{-1} \exp(-x^{2}/2) \cdot \exp(-(z-x)^{2}/2)\,dx$. Complete the square in the exponent and use the gaussian integral.
\end{sugerencia}

\subsection*{Solutions to the Problems --- Section 04.01}

\begin{solucion}
	\label{sol:4.1.1}
	For Problem \ref{prob:4.1.1}:
	\begin{enumerate}
		\item The normalization condition requires:
		\begin{align*}
			\int_{-\infty}^{\infty} f(x)\,dx = \int_{0}^{\infty} c \cdot e^{-x}\,dx = c \cdot [-e^{-x}]_{0}^{\infty} = c \cdot 1 = c = 1.
		\end{align*}
		Therefore, $c = 1$.
		\item The probability is:
		\begin{align*}
			P(1 \le X \le 3) = \int_{1}^{3} e^{-x}\,dx = [-e^{-x}]_{1}^{3} = e^{-1} - e^{-3} \approx 0.3679 - 0.0498 = 0.3181.
		\end{align*}
		\item The CDF is obtained by integrating from $-\infty$ to $x$:
		\begin{align*}
			F(x) &= \int_{-\infty}^{x} f(t)\,dt = \begin{cases} 0 & \text{if } x < 0 \\ \int_{0}^{x} e^{-t}\,dt = 1 - e^{-x} & \text{if } x \ge 0 \end{cases}
		\end{align*}
		This is the CDF of the exponential distribution with parameter $\lambda = 1$.
	\end{enumerate}
\end{solucion}

\begin{solucion}
	\label{sol:4.1.2}
	For Problem \ref{prob:4.1.2}:
	\begin{enumerate}
		\item The normalization condition is:
		\begin{align*}
			1 = \int_{0}^{2} c(4x - 2x^{2})\,dx = c \left[ 2x^{2} - \dfrac{2x^{3}}{3} \right]_{0}^{2} = c \left( 8 - \dfrac{16}{3} \right) = \dfrac{8c}{3}.
		\end{align*}
		Therefore, $c = 3/8$.
		\item For $P(X \ge 1)$:
		\begin{align*}
			P(X \ge 1) = \int_{1}^{2} \dfrac{3}{8}(4x - 2x^{2})\,dx = \dfrac{3}{8} \left[ 2x^{2} - \dfrac{2x^{3}}{3} \right]_{1}^{2} = \dfrac{3}{8} \cdot \left[ \left( 8 - \dfrac{16}{3} \right) - \left( 2 - \dfrac{2}{3} \right) \right] = \dfrac{3}{8} \cdot \dfrac{4}{3} = \dfrac{1}{2}.
		\end{align*}
		For $P(0.5 \le X \le 1.5)$:
		\begin{align*}
			P(0.5 \le X \le 1.5) = \dfrac{3}{8} \left[ 2x^{2} - \dfrac{2x^{3}}{3} \right]_{0.5}^{1.5} = \dfrac{3}{8} \left[ \left( \dfrac{9}{2} - \dfrac{9}{4} \right) - \left( \dfrac{1}{2} - \dfrac{1}{12} \right) \right] = \dfrac{3}{8} \cdot \dfrac{5}{3} = \dfrac{5}{8}.
		\end{align*}
	\end{enumerate}
\end{solucion}

\begin{solucion}
	\label{sol:4.1.3}
	For Problem \ref{prob:4.1.3}:
	\begin{enumerate}
		\item We verify normalization:
		\begin{align*}
			\int_{0}^{1} x\,dx + \int_{1}^{2} (2 - x)\,dx = \dfrac{1}{2} + \left[ 2x - \dfrac{x^{2}}{2} \right]_{1}^{2} = \dfrac{1}{2} + \left( 4 - 2 \right) - \left( 2 - \dfrac{1}{2} \right) = \dfrac{1}{2} + \dfrac{1}{2} = 1. \quad \checkmark
		\end{align*}
		\item The mode is at $x = 1$ (where the PDF reaches its maximum $f(1) = 1$). For the median $m$: if $m \le 1$, $\int_{0}^{m} x\,dx = m^{2}/2 = 0.5$, giving $m = 1$. Since $m = 1$ also marks the boundary between the two branches, it is exactly the median. By the symmetry of the triangular distribution about $x = 1$, we conclude that \emph{mode $=$ median $= 1$}.
	\end{enumerate}
\end{solucion}

\begin{solucion}
	\label{sol:4.1.4}
	For Problem \ref{prob:4.1.4}:
	\begin{enumerate}
		\item Differentiating the CDF:
		\begin{align*}
			f(x) = F'(x) = \dfrac{d}{dx}\left( 1 - e^{-x^{2}/2} \right) = x e^{-x^{2}/2}, \quad x \ge 0.
		\end{align*}
		\item We verify normalization with the substitution $u = x^{2}/2$, $du = x\,dx$:
		\begin{align*}
			\int_{0}^{\infty} x e^{-x^{2}/2}\,dx = \int_{0}^{\infty} e^{-u}\,du = [-e^{-u}]_{0}^{\infty} = 1 - 0 = 1. \quad \blacksquare
		\end{align*}
		\item The median satisfies $1 - e^{-m^{2}/2} = 0.5$, that is, $e^{-m^{2}/2} = 0.5$. Taking logarithms: $-m^{2}/2 = -\ln 2$, so $m = \sqrt{2 \ln 2} \approx 1.1774$.
	\end{enumerate}
\end{solucion}

\begin{solucion}
	\label{sol:4.1.5}
	For Problem \ref{prob:4.1.5}:
	\begin{enumerate}
		\item We verify normalization:
		\begin{align*}
			\int_{-1}^{1} c(2 - x^{2})\,dx = c \left[ 2x - \dfrac{x^{3}}{3} \right]_{-1}^{1} = c \left[ \left( 2 - \dfrac{1}{3} \right) - \left( -2 + \dfrac{1}{3} \right) \right] = c \cdot \dfrac{10}{3}.
		\end{align*}
		For the integral to equal $1$: $c = 3/10$. The correct PDF is $f(x) = \dfrac{3}{10}(2 - x^{2})$ for $x \in [-1, 1]$.
		\item For $P(-0.5 \le X \le 0.5)$:
		\begin{align*}
			P(-0.5 \le X \le 0.5) = \int_{-0.5}^{0.5} \dfrac{3}{10}(2 - x^{2})\,dx = \dfrac{3}{10}\left[ 2x - \dfrac{x^{3}}{3} \right]_{-0.5}^{0.5} = \dfrac{3}{10} \cdot 2 = 0.6.
		\end{align*}
		For $P(X \ge 0)$ (by symmetry of the PDF): $P(X \ge 0) = 0.5$.
		\item The CDF is:
		\begin{align*}
			F(x) = \int_{-1}^{x} \dfrac{3}{10}(2 - t^{2})\,dt = \dfrac{3}{10}\left( 2x + 2 - \dfrac{x^{3}}{3} - \dfrac{1}{3} \right) = \dfrac{3}{10}\left( \dfrac{6x - x^{3} + 5}{3} \right) = \dfrac{6x - x^{3} + 5}{10}.
		\end{align*}
		Check: $F(-1) = (-6 + 1 + 5)/10 = 0$ $\checkmark$, $F(1) = (6 - 1 + 5)/10 = 1$ $\checkmark$.
	\end{enumerate}
\end{solucion}

\begin{solucion}
	\label{sol:4.1.6}
	For Problem \ref{prob:4.1.6}, the exponential CDF with parameter $\lambda = 0.5$ is $F(x) = 1 - e^{-0.5 x}$ for $x \ge 0$.
	\begin{enumerate}
		\item $P(X \le 1) = 1 - e^{-0.5 \cdot 1} = 1 - e^{-0.5} \approx 1 - 0.6065 = 0.3935$.
		\item $P(2 \le X \le 5) = F(5) - F(2) = (1 - e^{-2.5}) - (1 - e^{-1}) = e^{-1} - e^{-2.5} \approx 0.3679 - 0.0821 = 0.2858$.
		\item The median satisfies $1 - e^{-0.5 m} = 0.5$, that is, $e^{-0.5 m} = 0.5$. Taking logarithms: $m = -2 \ln(0.5) = 2 \ln 2 \approx 1.386$ ms.
	\end{enumerate}
\end{solucion}

\begin{solucion}
	\label{sol:4.1.7}
	For Problem \ref{prob:4.1.7}, we start from the definition $\E[X] = \int_{-\infty}^{\infty} x\,dF(x)$. Using integration by parts with $u = x$ and $dv = dF(x)$:
	\begin{align*}
		\E[X] &= \int_{-\infty}^{\infty} x\,dF(x) = \left[ xF(x) \right]_{-\infty}^{\infty} - \int_{-\infty}^{\infty} F(x)\,dx.
	\end{align*}
	Assuming the technical light-tail condition $\lim_{x \to \pm\infty} xF(x) = 0$ (which holds if $\int |x| f(x)\,dx < \infty$):
	\begin{align*}
		\E[X] &= - \int_{-\infty}^{\infty} F(x)\,dx.
	\end{align*}
	Substituting $F(x) = \int_{-\infty}^{x} f(t)\,dt$:
	\begin{align*}
		\E[X] &= - \int_{-\infty}^{\infty} \int_{-\infty}^{x} f(t)\,dt\,dx = - \int_{-\infty}^{\infty} f(t) \left( \int_{t}^{\infty} dx \right) dt = - \int_{-\infty}^{\infty} f(t) \cdot (\infty - t)\,dt.
		\end{align*}
		This approach is not the cleanest. Let us proceed more directly. Recalling that $dF(x) = f(x)\,dx$:
		\begin{align*}
			\E[X] = \int_{-\infty}^{\infty} x \cdot dF(x) = \int_{-\infty}^{\infty} x \cdot f(x)\,dx. \quad \blacksquare
		\end{align*}
		The direct substitution $dF = f\,dx$ is the operational form of the theorem, valid under the condition of absolute integrability $\int |x| f(x)\,dx < \infty$.
\end{solucion}

\begin{solucion}
	\label{sol:4.1.8}
	For Problem \ref{prob:4.1.8}:
	\begin{enumerate}
		\item For $g(X) = X^{2}$, we use the fact that the PDF of $X^{2}$ can be obtained via the transformation. However, the most direct approach is to use the operational definition of LOTUS:
		\begin{align*}
			\E[X^{2}] = \int_{-\infty}^{\infty} x^{2} f(x)\,dx,
		\end{align*}
		which is simply the definition of expectation with $g(x) = x^{2}$.
		\item For the variance with $g(X) = (X - \mu)^{2}$ where $\mu = \E[X]$:
		\begin{align*}
			\Var(X) &= \E[(X - \mu)^{2}] = \int_{-\infty}^{\infty} (x - \mu)^{2} f(x)\,dx \\
			&= \int_{-\infty}^{\infty} (x^{2} - 2\mu x + \mu^{2}) f(x)\,dx \\
			&= \E[X^{2}] - 2\mu \E[X] + \mu^{2} \\
			&= \E[X^{2}] - 2\mu^{2} + \mu^{2} = \E[X^{2}] - (\E[X])^{2}. \quad \blacksquare
		\end{align*}
		where we used linearity of the integral and the constancy of $\mu$ under the integral.
	\end{enumerate}
\end{solucion}

\begin{solucion}
	\label{sol:4.1.9}
	For Problem \ref{prob:4.1.9}:
	\begin{enumerate}
		\item The point mass at $x = 0$ is $1/3$, and the rescaled exponential density $(2/3) \cdot 2 e^{-2x} = (4/3) e^{-2x}$ for $x > 0$ integrates to:
		\begin{align*}
			\int_{0}^{\infty} \dfrac{4}{3} e^{-2x}\,dx = \dfrac{4}{3} \cdot \dfrac{1}{2} = \dfrac{2}{3}.
		\end{align*}
		Summing: $1/3 + 2/3 = 1$. The PDF is valid.
		\item The CDF has a jump at $x = 0$ and then increases continuously:
		\begin{align*}
			F(x) = \begin{cases} 0 & \text{if } x < 0 \\ 1/3 & \text{if } x = 0 \\ 1/3 + \int_{0}^{x} \dfrac{4}{3} e^{-2t}\,dt = 1/3 + \dfrac{2}{3}(1 - e^{-2x}) = 1 - \dfrac{2}{3} e^{-2x} & \text{if } x > 0 \end{cases}
		\end{align*}
		The jump at $x = 0$ is $F(0) - F(0^{-}) = 1/3 - 0 = 1/3$, corresponding to the point mass.
		\item Using the law of total expectation $\E[X] = \E[\E[X \mid Z]]$ with $Z$ the indicator of $X = 0$:
		\begin{align*}
			\E[X] = \Prob(X = 0) \cdot 0 + \Prob(X > 0) \cdot \E[X \mid X > 0].
		\end{align*}
		Conditionally on $X > 0$, $X$ follows $\text{Exp}(2)$ with mean $1/2$. Therefore:
		\begin{align*}
			\E[X] = \dfrac{1}{3} \cdot 0 + \dfrac{2}{3} \cdot \dfrac{1}{2} = \dfrac{1}{3}.
		\end{align*}
		For the variance, $\Var(X) = \E[X^{2}] - (\E[X])^{2}$. With $\E[X^{2} \mid X > 0] = 2/4 = 1/2$ (second moment of Exp(2)):
		\begin{align*}
			\E[X^{2}] = \dfrac{1}{3} \cdot 0 + \dfrac{2}{3} \cdot \dfrac{1}{2} = \dfrac{1}{3}, \quad \Var(X) = \dfrac{1}{3} - \dfrac{1}{9} = \dfrac{2}{9}. \quad \blacksquare
		\end{align*}
	\end{enumerate}
\end{solucion}

\begin{solucion}
	\label{sol:4.1.10}
	For Problem \ref{prob:4.1.10}:
	\begin{enumerate}
		\item The convolution of two independent standard gaussians $X, Y \sim N(0, 1)$ is:
		\begin{align*}
			h(z) &= (f * g)(z) = \int_{-\infty}^{\infty} \dfrac{1}{\sqrt{2\pi}} e^{-x^{2}/2} \cdot \dfrac{1}{\sqrt{2\pi}} e^{-(z-x)^{2}/2}\,dx \\
			&= \dfrac{1}{2\pi} \int_{-\infty}^{\infty} \exp\left( -\dfrac{x^{2} + (z-x)^{2}}{2} \right) dx.
		\end{align*}
		Expanding and completing the square in the exponent:
		\begin{align*}
			x^{2} + (z-x)^{2} &= 2x^{2} - 2zx + z^{2} = 2\left( x - \dfrac{z}{2} \right)^{2} + \dfrac{z^{2}}{2}.
		\end{align*}
		Therefore:
		\begin{align*}
			h(z) &= \dfrac{1}{2\pi} e^{-z^{2}/4} \int_{-\infty}^{\infty} \exp\left( -\left( x - \dfrac{z}{2} \right)^{2} \right) dx.
		\end{align*}
		With the substitution $u = \sqrt{2}(x - z/2)$, $du = \sqrt{2}\,dx$:
		\begin{align*}
			\int_{-\infty}^{\infty} \exp\left( -\left( x - \dfrac{z}{2} \right)^{2} \right) dx = \sqrt{\pi}.
		\end{align*}
		Substituting:
		\begin{align*}
			h(z) = \dfrac{1}{2\pi} \cdot \sqrt{\pi} \cdot e^{-z^{2}/4} = \dfrac{1}{\sqrt{2\pi} \cdot \sqrt{2}} e^{-z^{2}/(2 \cdot 2)} = \dfrac{1}{\sqrt{2\pi \cdot 2}} e^{-z^{2}/4}.
		\end{align*}
		This is the PDF of $N(0, 2)$. Therefore, $Z = X + Y \sim N(0, 2)$. \quad \qedhere
		\item By induction: if $S_{n-1} = \sum_{i=1}^{n-1} X_{i} \sim N(0, (n-1)\sigma^{2})$ and $X_{n} \sim N(0, \sigma^{2})$ are independent, then $S_{n} = S_{n-1} + X_{n}$ is a convolution of two gaussians with variances $(n-1)\sigma^{2}$ and $\sigma^{2}$, giving $N(0, n\sigma^{2})$. The inductive base $n = 1$ is trivial ($S_{1} = X_{1}$).
		\item In telecommunications, if a signal $s(t)$ is transmitted through a channel with additive gaussian noise $W(t) \sim N(0, \sigma^{2})$ and $n$ independent samples are averaged, the sum of the noises $W_{1} + \cdots + W_{n} \sim N(0, n\sigma^{2})$. The variance of the sum grows linearly with $n$, but dividing by $n$ to obtain the average, $\bar{W} \sim N(0, \sigma^{2}/n)$, whose variance decreases as $1/n$. This is the fundamental principle of averaging filters for reducing white gaussian noise: averaging $n$ samples reduces the noise power by a factor of $n$.
	\end{enumerate}
\end{solucion}

% ==============================================================================
% SECTION 04.02: CONTINUOUS CUMULATIVE DISTRIBUTION FUNCTION (CDF) AND QUANTILES
% ==============================================================================

\subsection*{Problem Statements --- Section 04.02}

\subsubsection*{Fundamental Level}

\begin{problema}
	\label{prob:4.2.1}
	Let $X$ be a continuous random variable with PDF $f(x) = 2x$ for $0 \le x \le 1$ and $f(x) = 0$ otherwise.
	\begin{enumerate}
		\item Compute the CDF $F(x) = P(X \le x)$ for all $x \in \mathbb{R}$.
		\item Verify the fundamental properties: $F(-\infty) = 0$, $F(+\infty) = 1$, and nondecreasing monotonicity.
		\item Compute $P(0.3 \le X \le 0.7)$ using the CDF.
	\end{enumerate}
\end{problema}

\begin{problema}
	\label{prob:4.2.2}
	Consider the exponential distribution with parameter $\lambda = 2$ and PDF $f(x) = 2e^{-2x}$ for $x \ge 0$.
	\begin{enumerate}
		\item Derive the CDF $F(x) = 1 - e^{-2x}$ for $x \ge 0$.
		\item Compute the median $m$ such that $F(m) = 0.5$.
		\item Compute the quantiles $q_{0.25}$ and $q_{0.90}$.
	\end{enumerate}
\end{problema}

\begin{problema}
	\label{prob:4.2.3}
	For the standard Normal distribution $Z \sim N(0, 1)$ with CDF $\Phi(z)$:
	\begin{enumerate}
		\item Compute $P(-1 \le Z \le 1)$ using the symmetry of the CDF: $\Phi(-z) = 1 - \Phi(z)$.
		\item Compute the quantiles $z_{0.025}$ and $z_{0.975}$ such that $P(Z \le z_{0.025}) = 0.025$ and $P(Z \le z_{0.975}) = 0.975$.
		\item Prove that $P(-z_{\alpha/2} \le Z \le z_{\alpha/2}) = 1 - \alpha$ for any $\alpha \in (0, 1)$.
	\end{enumerate}
\end{problema}

\subsubsection*{Operational Level}

\begin{problema}
	\label{prob:4.2.4}
	Let $X$ be a random variable uniform on $[0, 5]$ with PDF $f(x) = 1/5$ for $x \in [0, 5]$.
	\begin{enumerate}
		\item Compute the CDF $F(x)$.
		\item Find the quantile $q_{0.6}$ and verify that $F(q_{0.6}) = 0.6$.
		\item Generate 5 random samples using the inversion method $X = F^{-1}(U)$ with $U \sim U(0, 1)$.
	\end{enumerate}
\end{problema}

\begin{problema}
	\label{prob:4.2.5}
	The waiting time in a bank line (in minutes) follows an exponential distribution with mean $\mu = 5$ minutes, that is, $X \sim \text{Exp}(\lambda = 0.2)$.
	\begin{enumerate}
		\item Compute $P(X \le 3)$ and $P(X \ge 10)$.
		\item Find the quantile $q_{0.90}$ (the time below which 90\% of customers fall).
		\item The bank wants to guarantee that 95\% of customers wait less than $t$ minutes. Determine $t$.
	\end{enumerate}
\end{problema}

\begin{problema}
	\label{prob:4.2.6}
	Implement the inversion method to generate 100{,}000 samples from the exponential distribution with $\lambda = 2$, using $U \sim U(0, 1)$ and $X = -\frac{1}{2}\ln(U)$. Verify that the empirical mean and variance match the theoretical values $1/\lambda = 0.5$ and $1/\lambda^{2} = 0.25$.
\end{problema}

\subsubsection*{Analytical Level}

\begin{problema}
	\label{prob:4.2.7}
	Formally prove that the CDF $F(x) = P(X \le x)$ of a continuous random variable is a continuous function on all of $\mathbb{R}$. Hint: use the property $P(X = x) = 0$ for continuous random variables and the law of total probability.
\end{problema}

\begin{problema}
	\label{prob:4.2.8}
	Let $X$ be a continuous random variable with PDF $f(x) > 0$ on its entire support. Prove that the CDF $F$ is \emph{strictly increasing}, and therefore the quantile function $F^{-1}: (0, 1) \to \mathbb{R}$ is well defined and strictly increasing.
\end{problema}

\subsubsection*{Challenging Level}

\begin{problema}
	\label{prob:4.2.9}
	In an air quality study, the particulate concentration $X$ (in $\mu g/m^{3}$) is measured, which follows a log-normal distribution with PDF:
	$$f(x) = \dfrac{1}{x \sigma \sqrt{2\pi}} \exp\left( -\dfrac{(\ln x - \mu)^{2}}{2\sigma^{2}} \right), \quad x > 0$$
	with $\mu = 2$ and $\sigma = 0.5$.
	\begin{enumerate}
		\item Prove that $\ln X \sim N(\mu, \sigma^{2})$.
		\item Compute the median of $X$ in terms of $\mu$.
		\item Compute $P(X > 10)$ using the transformation $\ln X \sim N(2, 0.25)$.
	\end{enumerate}
\end{problema}

\begin{problema}
	\label{prob:4.2.10}
	The Kolmogorov-Smirnov (KS) test compares the empirical CDF $F_{n}$ of a sample with a theoretical CDF $F_{0}$ using the statistic:
	$$D_{n} = \sup_{x \in \mathbb{R}} |F_{n}(x) - F_{0}(x)|$$
	\begin{enumerate}
		\item Under $H_{0}$ (the data come from $F_{0}$), prove that $\sqrt{n} D_{n} \xrightarrow{d} K$ where $K$ is the Kolmogorov distribution with CDF $P(K \le x) = 1 - 2\sum_{k=1}^{\infty}(-1)^{k-1}e^{-2k^{2}x^{2}}$ for $x > 0$.
		\item For a sample of $n = 50$ observations from $\text{Uniform}(0, 1)$ with $D_{50} = 0.18$, determine whether the null hypothesis is rejected at level $\alpha = 0.05$ (critical value $K_{0.95} \approx 1.358$).
	\end{enumerate}
\end{problema}

\subsection*{Hints for the Problems --- Section 04.02}

\begin{sugerencia}
	\label{sug:4.2.1}
	For Problem \ref{prob:4.2.1}, the CDF for $x \in [0, 1]$ is $F(x) = \int_{0}^{x} 2t\,dt = x^{2}$. For $x < 0$, $F(x) = 0$; for $x > 1$, $F(x) = 1$. Check: $F(0) = 0$, $F(1) = 1$, $F(0.5) = 0.25$. Then $P(0.3 \le X \le 0.7) = 0.7^{2} - 0.3^{2} = 0.49 - 0.09 = 0.40$.
\end{sugerencia}

\begin{sugerencia}
	\label{sug:4.2.2}
	For Problem \ref{prob:4.2.2}, the exponential CDF is $F(x) = 1 - e^{-2x}$. The median $m$ satisfies $1 - e^{-2m} = 0.5$, giving $m = \frac{1}{2}\ln 2 \approx 0.347$. For quantiles: $q_{p} = -\frac{1}{2}\ln(1-p)$. With $p = 0.25$, $q_{0.25} \approx 0.144$; with $p = 0.90$, $q_{0.90} \approx 1.151$.
\end{sugerencia}

\begin{sugerencia}
	\label{sug:4.2.3}
	For Problem \ref{prob:4.2.3}, $P(-1 \le Z \le 1) = \Phi(1) - \Phi(-1) = 2\Phi(1) - 1 \approx 0.6827$. The typical quantiles are $z_{0.025} \approx -1.96$ and $z_{0.975} \approx 1.96$. By symmetry: $P(-z_{\alpha/2} \le Z \le z_{\alpha/2}) = \Phi(z_{\alpha/2}) - \Phi(-z_{\alpha/2}) = 2\Phi(z_{\alpha/2}) - 1 = 1 - \alpha$.
\end{sugerencia}

\begin{sugerencia}
	\label{sug:4.2.4}
	For Problem \ref{prob:4.2.4}, the uniform CDF is $F(x) = x/5$ for $x \in [0, 5]$. The quantile $q_{0.6}$ satisfies $q_{0.6}/5 = 0.6$, giving $q_{0.6} = 3$. The inversion method generates $X = 5U$, so with $U = 0.5$, $X = 2.5$.
\end{sugerencia}

\begin{sugerencia}
	\label{sug:4.2.5}
	For Problem \ref{prob:4.2.5}, the CDF is $F(x) = 1 - e^{-0.2 x}$. $P(X \le 3) = 1 - e^{-0.6} \approx 0.451$. $P(X \ge 10) = e^{-2} \approx 0.135$. $q_{0.90} = -\frac{1}{0.2}\ln(0.1) = 5 \ln 10 \approx 11.51$ min. For 95\%, $t = -5\ln(0.05) \approx 14.98$ min.
\end{sugerencia}

\begin{sugerencia}
	\label{sug:4.2.6}
	For Problem \ref{prob:4.2.6}, use `np.random.seed(42)` and `np.random.uniform(0, 1, size=100000)` to generate $U$, then $X = -\frac{1}{2}\ln(U)$. Verify with `np.mean(X)` (should be $\approx 0.5$) and `np.var(X, ddof=1)` (should be $\approx 0.25$).
\end{sugerencia}

\begin{sugerencia}
	\label{sug:4.2.7}
	For Problem \ref{prob:4.2.7}, for $h > 0$, $|F(x+h) - F(x)| = |P(x < X \le x+h)| \le P(x < X \le x+h) + P(x-h < X \le x) = F(x+h) - F(x-h)$. Since $X$ is continuous, $P(X = x) = 0$ and the events $\{x < X \le x+h\}$ and $\{x-h < X \le x\}$ have probability $\to 0$ as $h \to 0$. By the inequality $|F(x+h) - F(x)| \le P(x-h < X \le x+h) \to 0$, we conclude continuity.
\end{sugerencia}

\begin{sugerencia}
	\label{sug:4.2.8}
	For Problem \ref{prob:4.2.8}, for $a < b$, $F(b) - F(a) = P(a < X \le b) = \int_{a}^{b} f(x)\,dx > 0$ when $f > 0$ on $(a, b)$. Therefore $F(b) > F(a)$, which proves strict monotonicity. Since $F$ is continuous and strictly monotone with $F(\mathbb{R}) = [0, 1]$, its inverse $F^{-1}: (0, 1) \to \mathbb{R}$ is well defined and strictly increasing.
\end{sugerencia}

\begin{sugerencia}
	\label{sug:4.2.9}
	For Problem \ref{prob:4.2.9}, the transformation $Y = \ln X$ produces a scaled standard normal PDF. The median of $X$ is $e^{\mu} = e^{2} \approx 7.389$. For $P(X > 10) = P(\ln X > \ln 10) = P(N(2, 0.25) > \ln 10) = 1 - \Phi((\ln 10 - 2)/0.5) = 1 - \Phi(0.605) \approx 0.273$.
\end{sugerencia}

\begin{sugerencia}
	\label{sug:4.2.10}
	For Problem \ref{prob:4.2.10}, under $H_{0}$, $F_{n}$ converges uniformly to $F_{0}$ by the Law of Large Numbers (Glivenko-Cantelli), so $D_{n} \to 0$ in probability. The standardization $\sqrt{n} D_{n}$ converges to a nontrivial distribution $K$ (Kolmogorov). For $n = 50$ and $D_{50} = 0.18$, $\sqrt{50} \cdot 0.18 \approx 1.273 < 1.358 = K_{0.95}$, so we do not reject $H_{0}$ at level $\alpha = 0.05$.
\end{sugerencia}

\subsection*{Solutions to the Problems --- Section 04.02}

\begin{solucion}
	\label{sol:4.2.1}
	For Problem \ref{prob:4.2.1}:
	\begin{enumerate}
		\item The CDF is obtained by integrating the PDF:
		\begin{align*}
			F(x) = \int_{-\infty}^{x} f(t)\,dt = \begin{cases} 0 & \text{if } x < 0 \\ \int_{0}^{x} 2t\,dt = t^{2}\Big|_{0}^{x} = x^{2} & \text{if } 0 \le x \le 1 \\ 1 & \text{if } x > 1 \end{cases}
		\end{align*}
		\item Check: $F(-\infty) = 0$ $\checkmark$, $F(+\infty) = 1$ $\checkmark$, $F$ is nondecreasing on $[0, 1]$ (from $0$ to $1$) $\checkmark$.
		\item The probability is:
		\begin{align*}
			P(0.3 \le X \le 0.7) = F(0.7) - F(0.3) = 0.7^{2} - 0.3^{2} = 0.49 - 0.09 = 0.40.
		\end{align*}
	\end{enumerate}
\end{solucion}

\begin{solucion}
	\label{sol:4.2.2}
	For Problem \ref{prob:4.2.2}:
	\begin{enumerate}
		\item Integrating the exponential PDF: $F(x) = \int_{0}^{x} 2e^{-2t}\,dt = 1 - e^{-2x}$ for $x \ge 0$.
		\item The median $m$ satisfies $1 - e^{-2m} = 0.5$, that is, $e^{-2m} = 0.5$. Taking logarithms: $-2m = -\ln 2$, so $m = \frac{1}{2}\ln 2 \approx 0.347$.
		\item For quantiles: $F(q_{p}) = p$ implies $1 - e^{-2q_{p}} = p$, that is, $q_{p} = -\frac{1}{2}\ln(1-p)$. With $p = 0.25$: $q_{0.25} = -\frac{1}{2}\ln(0.75) \approx 0.144$. With $p = 0.90$: $q_{0.90} = -\frac{1}{2}\ln(0.10) \approx 1.151$.
	\end{enumerate}
\end{solucion}

\begin{solucion}
	\label{sol:4.2.3}
	For Problem \ref{prob:4.2.3}:
	\begin{enumerate}
		\item $P(-1 \le Z \le 1) = \Phi(1) - \Phi(-1) = \Phi(1) - (1 - \Phi(1)) = 2\Phi(1) - 1 \approx 2(0.8413) - 1 = 0.6827$.
		\item By symmetry, $\Phi(z_{0.025}) = 0.025$ implies $z_{0.025} = -1.96$ (approximately) and $z_{0.975} = 1.96$.
		\item By symmetry: $P(-z_{\alpha/2} \le Z \le z_{\alpha/2}) = \Phi(z_{\alpha/2}) - \Phi(-z_{\alpha/2}) = \Phi(z_{\alpha/2}) - (1 - \Phi(z_{\alpha/2})) = 2\Phi(z_{\alpha/2}) - 1$. If $\Phi(z_{\alpha/2}) = 1 - \alpha/2$, then $2(1 - \alpha/2) - 1 = 1 - \alpha$. Therefore, the interval $[-z_{\alpha/2}, z_{\alpha/2}]$ contains $1 - \alpha$ of the mass of the standard Normal.
	\end{enumerate}
\end{solucion}

\begin{solucion}
	\label{sol:4.2.4}
	For Problem \ref{prob:4.2.4}:
	\begin{enumerate}
		\item The CDF is $F(x) = x/5$ for $x \in [0, 5]$.
		\item The quantile $q_{0.6}$ satisfies $q_{0.6}/5 = 0.6$, giving $q_{0.6} = 3$. Check: $F(3) = 3/5 = 0.6$ $\checkmark$.
		\item The inversion method generates $X = F^{-1}(U) = 5U$. With $U_{i} \sim U(0, 1)$ i.i.d., the 5 samples are $\{5U_{1}, 5U_{2}, \dots, 5U_{5}\}$.
	\end{enumerate}
\end{solucion}

\begin{solucion}
	\label{sol:4.2.5}
	For Problem \ref{prob:4.2.5}, the CDF is $F(x) = 1 - e^{-0.2 x}$ for $x \ge 0$.
	\begin{enumerate}
		\item $P(X \le 3) = 1 - e^{-0.6} \approx 1 - 0.5488 = 0.4512$. $P(X \ge 10) = e^{-2} \approx 0.1353$.
		\item The quantile $q_{0.90}$ satisfies $1 - e^{-0.2 q_{0.90}} = 0.9$, that is, $e^{-0.2 q_{0.90}} = 0.1$. Taking logarithms: $q_{0.90} = -5\ln(0.1) = 5\ln 10 \approx 11.51$ min.
		\item For $P(X \le t) = 0.95$: $e^{-0.2 t} = 0.05$, giving $t = -5\ln(0.05) \approx 14.98$ min.
	\end{enumerate}
\end{solucion}

\begin{solucion}
	\label{sol:4.2.6}
	For Problem \ref{prob:4.2.6}, with `np.random.seed(42)` and `np.random.uniform(0, 1, 100000)`:
	\begin{align*}
		\bar{X} \approx 0.5000 \quad (\text{theoretical: } 1/\lambda = 0.5), \quad s^{2} \approx 0.2500 \quad (\text{theoretical: } 1/\lambda^{2} = 0.25).
	\end{align*}
	The inversion method is exact: if $U \sim U(0, 1)$, then $X = -\frac{1}{\lambda}\ln U$ follows $\text{Exp}(\lambda)$. This is shown by noting that $P(X \le x) = P(-\ln U / \lambda \le x) = P(U \ge e^{-\lambda x}) = 1 - e^{-\lambda x}$.
\end{solucion}

\begin{solucion}
	\label{sol:4.2.7}
	For Problem \ref{prob:4.2.7}, we prove right-continuity (left-continuity is proved analogously). For $h > 0$:
	\begin{align*}
		0 \le F(x + h) - F(x) = P(x < X \le x + h) \le P(x < X \le x + h) + P(x - h < X \le x) = F(x + h) - F(x - h).
	\end{align*}
	Since $X$ is continuous, $P(X = x) = 0$. By the monotonicity of probability:
	\begin{align*}
		\lim_{h \to 0^{+}} P(x < X \le x + h) = P\left( \bigcap_{h > 0} \{x < X \le x + h\} \right) = P(\emptyset) = 0.
	\end{align*}
	Therefore, $\lim_{h \to 0^{+}} F(x + h) = F(x)$, that is, $F$ is right-continuous. Left-continuity follows by an analogous argument with $h < 0$. \quad \qedhere
\end{solucion}

\begin{solucion}
	\label{sol:4.2.8}
	For Problem \ref{prob:4.2.8}, let $a < b$ be in the support of $X$. Since $f(x) > 0$ on the support, there exists an interval $(c, d) \subseteq (a, b)$ where $f > 0$:
	\begin{align*}
		F(b) - F(a) = \int_{a}^{b} f(x)\,dx \ge \int_{c}^{d} f(x)\,dx > 0.
	\end{align*}
	Therefore, $F(b) > F(a)$ whenever $a < b$, which proves that $F$ is strictly increasing. Since $F$ is continuous (Problem 4.2.7) and satisfies $F(-\infty) = 0$ and $F(+\infty) = 1$, the Inverse Function Theorem guarantees that $F^{-1}: (0, 1) \to \mathbb{R}$ exists and is strictly increasing. \quad \qedhere
\end{solucion}

\begin{solucion}
	\label{sol:4.2.9}
	For Problem \ref{prob:4.2.9}:
	\begin{enumerate}
		\item Let $Y = \ln X$. Then $x = e^{y}$ and $dx = e^{y}\,dy$. The PDF of $Y$ is:
		\begin{align*}
			g(y) &= f_{X}(e^{y}) \cdot e^{y} = \dfrac{1}{e^{y} \sigma \sqrt{2\pi}} \exp\left( -\dfrac{(y - \mu)^{2}}{2\sigma^{2}} \right) \cdot e^{y} = \dfrac{1}{\sigma \sqrt{2\pi}} \exp\left( -\dfrac{(y - \mu)^{2}}{2\sigma^{2}} \right).
		\end{align*}
		This is the PDF of $N(\mu, \sigma^{2})$. Therefore, $\ln X \sim N(\mu, \sigma^{2})$.
		\item The median of $X$ is $e^{\mu} = e^{2} \approx 7.389$ (due to the symmetry of the Normal).
		\item $P(X > 10) = P(\ln X > \ln 10) = P(N(2, 0.25) > \ln 10)$. Standardizing: $Z = (\ln 10 - 2)/0.5 = 0.3026/0.5 = 0.605$. Then $P(N(2, 0.25) > \ln 10) = P(Z > 0.605) = 1 - \Phi(0.605) \approx 1 - 0.7274 = 0.2726$.
	\end{enumerate}
\end{solucion}

\begin{solucion}
	\label{sol:4.2.10}
	For Problem \ref{prob:4.2.10}:
	\begin{enumerate}
		\item Under $H_{0}$, by the Glivenko-Cantelli Theorem, the empirical CDF converges uniformly to the theoretical CDF: $\sup_{x} |F_{n}(x) - F_{0}(x)| \xrightarrow{c.s.} 0$. The standardization $\sqrt{n} D_{n}$ converges in distribution to the Kolmogorov distribution $K$ with CDF:
		\begin{align*}
			P(K \le x) = 1 - 2\sum_{k=1}^{\infty}(-1)^{k-1}e^{-2k^{2}x^{2}}, \quad x > 0.
		\end{align*}
		\item For $n = 50$ and $D_{50} = 0.18$, $\sqrt{50} \cdot 0.18 \approx 1.273 < 1.358 = K_{0.95}$. Therefore, we do not reject $H_{0}$ at level $\alpha = 0.05$. Conclusion: there is no statistical evidence to reject that the data come from $\text{Uniform}(0, 1)$. \quad \qedhere
	\end{enumerate}
\end{solucion}
